% \begin{figure*}[t]
\centering
\begin{tikzpicture}[x=1cm,y=1cm,>=Stealth,
 box/.style={draw=methodblue,rounded corners=1pt,align=center,minimum height=0.85cm,font=\normalsize},
 smallbox/.style={box,text width=3.5cm},
 arr/.style={->,thick,draw=methodblue}]
\node[smallbox] (dense) at (1.8,1.0) {Dense image input\\$(c,y,x,\text{threshold})$};
\node[smallbox] (conv) at (1.8,-0.25) {Convolutional channels\\retain spatial sampling};
\node[box,text width=3.25cm] (semantic) at (6.0,1.0) {Semantic pairing\\first layer / classifier input};
\node[box,text width=3.25cm] (stages) at (6.0,-0.25) {Multiscale matchings\\deeper layers};
\node[box,text width=3.4cm] (fixed) at (10.4,0.4) {Fixed predecessor indices\\train gate functions};
\node[box,text width=2.4cm] (hard) at (14.1,0.4) {Harden gates\\Boolean inference};
\draw[arr] (dense) -- (semantic);
\draw[arr] (conv.east) -- ++(0.35,0) |- (semantic.west);
\draw[arr] (semantic) -- (stages);
\draw[arr] (semantic.east) -- ++(0.4,0) |- (fixed.west);
\draw[arr] (stages.east) -- ++(0.4,0) |- (fixed.west);
\draw[arr] (fixed) -- (hard);
\node[align=center,font=\normalsize,text=controlgray] at (7.6,-1.2)
 {Stage priority: degree spread $\rightarrow$ ancestry novelty $\rightarrow$ nominal scale};
\end{tikzpicture}
\caption{Overview of LogicConnect. Semantic pairing constructs input connections, and multiscale matching organizes deeper connections. The convolutional construction preserves spatial sampling and shared logic trees. Training optimizes gate functions on the resulting fixed graph.}
\label{fig:overview}
\end{figure*}

\section{Methodology}
\label{sec:method}

\textcolor{red}{LogicConnect constructs the connectivity of a DLGN \emph{before training},
with the goal of exposing each gate to more informative combinations of
signals while retaining the simplicity of fixed connectivity}. As illustrated
in Figure~\ref{fig:overview}, LogicConnect takes the encoded input
representation and the dimensions of each network layer and returns the two
fixed predecessor indices required by every gate. Once these connections are
generated, the graph is frozen and training proceeds using the original DLGN
gate parameterization. Therefore, LogicConnect does not introduce additional
routing parameters, change the number of gates, or modify the gate-training
procedure.

The key challenge is that the information available for connectivity design
changes with network depth. At the network input, signals retain explicit
relationships inherited from the data representation, such as their spatial
location, channel, and encoding. LogicConnect exploits these relationships
through \emph{semantic input pairing}, avoiding an arbitrary selection of inputs. After signals pass through logic gates, however, this direct
interpretation is progressively lost. Rather than imposing the same semantic
rule at every layer, LogicConnect organizes subsequent connections through
\emph{degree-balanced multiscale matchings}. These matchings systematically
combine predecessors while controlling their reuse based on structural ancestry. 

Importantly, the construction is entirely architecture-driven. LogicConnect
does not observe class labels, training loss, validation or test accuracy, or
trained gate parameters when selecting connections. It also does not search
over individual edges using accuracy feedback or perform post-hoc edge
swapping. Instead, it selects from predefined structured pairing and matching
stages using deterministic structural criteria. Thus, the one-time
construction is separated from learning: connectivity is designed offline,
while conventional DLGN training subsequently determines only the Boolean
functions implemented by the fixed graph. 
% The following subsections describe the graph representation, semantic input construction, multiscale matching, and its application to convolutional DLGNs

\subsection{Fixed-Graph Representation}
\label{sec:representation}

We first define the fixed connectivity used throughout the network. For a
layer $l$ with $m_l$ two-input gates and $n_l$ available predecessors, we
represent its connections by
$I_l\in\{0,\ldots,n_l-1\}^{2\times m_l}$. For every gate $j$, the two
entries $I_l[0,j]$ and $I_l[1,j]$ identify distinct predecessor signals.
Following the original DLGN formulation~\cite{petersen2022}, the output of
gate $j$ during training is computed as
%
\begin{equation}
h_{l,j} =
\sum_{q=0}^{15}
\mathrm{softmax}(\theta_{l,j})_q\,
\widetilde{f}_q
\!\left(
h_{l-1,I_l[0,j]},
h_{l-1,I_l[1,j]}
\right),
\label{eq:gate}
\end{equation}
%
where $\theta_{l,j}$ contains the trainable logits over the 16 two-input
Boolean functions and $\widetilde{f}_q$ denotes their continuous
relaxation. After training, the function with the highest weight is selected
for discrete inference. LogicConnect modifies only the fixed predecessor
indices $I_l$; the gate parameterization, optimization objective,
discretization procedure, and output aggregation remain unchanged. This
allows us to isolate the effect of connectivity from changes to the DLGN
training formulation.

To preserve information about how signals are related as they propagate
through the fixed graph, we associate every predecessor with a
\emph{structural ancestry set}. At the network input, each encoded signal is
assigned a source identity describing the raw value from which it originates.
For image inputs, for example, multiple thermometer-encoded thresholds of
the same pixel and channel share the same source identity. The ancestry of an
input therefore contains its source label. For a gate receiving predecessors
$a$ and $b$, its ancestry is propagated simply as
%
\begin{equation}
A_{l,j}=A_{l-1,a}\cup A_{l-1,b}.
\label{eq:ancestry}
\end{equation}
%
Importantly, $A_{l,j}$ describes the set of sources that are
\emph{structurally reachable} by the gate, rather than the inputs that
functionally determine its output. The learned Boolean function may suppress
or ignore one of its predecessors. We use this structural information only
during offline connectivity construction; it does not introduce trainable
parameters or alter the gate computation.

\subsection{Semantic Input Pairing}
\label{sec:semantic}

The first layer is the only part of the network where the identity of each
signal is known directly from the input representation. Rather than treating
these inputs as arbitrary bit indices, we use their source metadata to expose
the gates to relationships already present in the encoded image. We denote an
encoded input by $x_{p,c,t}$, where $p$ identifies the spatial position, $c$
the channel, and $t$ the encoding threshold. All thresholds originating from
the same raw pixel--channel value share the same source identity,
$\sigma(x_{p,c,t})=(p,c)$. Therefore, pairs satisfying
$\sigma(x_i)=\sigma(x_j)$ are excluded, since they originate from the same
underlying image value and provide limited complementary information.

Using this metadata, we construct a collection of semantic pairing stages.
Spatial stages connect inputs whose positions differ by power-of-two
displacements,
%
\begin{equation}
\Delta = \{1,2,4,\ldots\},
\label{eq:semantic_stride}
\end{equation}
%
up to half of the corresponding image dimension. Horizontal and vertical
displacements are considered separately, with cyclic boundaries used to avoid
systematically favoring particular image regions. Additional stages capture
cross-channel and cross-threshold relationships, including aligned channels
and threshold combinations with spatial displacement. Hence, the candidate
set can be viewed as
%
\begin{equation}
\mathcal{C}
=
\mathcal{C}_{\mathrm{spatial}}
\cup
\mathcal{C}_{\mathrm{channel}}
\cup
\mathcal{C}_{\mathrm{threshold}},
\label{eq:semantic_candidates}
\end{equation}
%
where every pair in $\mathcal{C}$ satisfies the distinct-source constraint
above. This exposes the first-layer gates to multiple complementary notions
of input structure instead of relying on a single local pairing rule.

Semantic structure alone, however, does not guarantee balanced predecessor
usage. Let $d_i$ denote the number of times predecessor $i$ has already been
selected. For a candidate pair $(i,j)$, LogicConnect prioritizes smaller
%
\begin{equation}
d_i+d_j,
\label{eq:pair_degree}
\end{equation}
%
followed by the smaller value of $\max(d_i,d_j)$, and finally the pair indices
to resolve remaining ties. This degree-aware ordering prevents a small subset
of semantically favorable inputs from dominating the first layer. Candidate
stages are visited cyclically so that spatial, channel, and encoding
relationships all contribute to the final graph.

To keep the construction scalable, we avoid searching the complete
$\mathcal{O}(n^2)$ space of possible predecessor pairs. At each selection,
the constructor examines at most 64 unused candidates following a fixed
seeded affine traversal and applies the degree criteria above. Once the
distinct pairs of a stage are exhausted, repeated pairs are allowed under the
same ordering. Importantly, this bounded search is purely structural:
class labels, training loss, validation accuracy, and trained gate parameters
are never used to select a pair. The first-layer connectivity is therefore
fully determined before training from the input metadata and the requested
network dimensions.


\subsection{Degree-Balanced Multiscale Matching}
\label{sec:multiscale}

Semantic pairing is well-defined at the network input because every signal
retains explicit spatial, channel, and encoding metadata. This interpretation
does not persist in deeper layers, where each predecessor represents the
output of several previously combined logic gates. Reapplying the same
input-level semantic rules would therefore impose relationships that are no
longer directly supported by the signal representation. Instead, LogicConnect
uses a regular \emph{multiscale matching} construction that progressively
mixes hidden signals while controlling how frequently each predecessor is
reused.

For a layer with $n$ available predecessors, we consider
$K=\lceil\log_2 n\rceil$ nominal interaction scales. When $n$ is a power of
two, scale $s$ forms the matching
%
\begin{equation}
\mathcal{M}_s =
\left\{(i,i\oplus 2^s)\right\},
\qquad
s=0,\ldots,\log_2 n-1,
\label{eq:xor_matching}
\end{equation}
%
where each unordered pair is retained only once. Consequently, every stage
contains $n/2$ disjoint pairs and every predecessor appears exactly once.
Different values of $s$ expose the gates to progressively different index
separations. For example, with eight predecessors,
%
\begin{align*}
s=0 &: (0,1),(2,3),(4,5),(6,7),\\
s=1 &: (0,2),(1,3),(4,6),(5,7),\\
s=2 &: (0,4),(1,5),(2,6),(3,7).
\end{align*}
%
This structure provides regular mixing at multiple scales without searching
over arbitrary edge combinations.

For widths that are not powers of two, we preserve the same principle using
affine traversals of the predecessor indices. A traversal follows
$\pi(k)=(o+kr)\bmod n$, where the step $r$ is chosen to be coprime with $n$,
and adjacent elements of the traversal are paired. This guarantees that the
traversal visits every predecessor before repeating. For odd $n$, one
predecessor necessarily remains unmatched in each complete stage; the
unmatched position is rotated across stages to avoid systematically excluding
the same signal. Thus, a complete matching contains $\lfloor n/2\rfloor$
pairs. If fewer gates remain than required for a complete stage, LogicConnect
uses a final partial matching and selects disjoint pairs according to endpoint
degree sum, maximum endpoint degree, and index order.

The remaining question is \emph{which matching stage should be used next}.
LogicConnect follows a strict selection hierarchy. First, it prioritizes
balanced predecessor reuse. Let $d_i$ denote the current number of times
predecessor $i$ has been selected, and let $\delta_i(\mathcal{M}_s)$ denote
its additional occurrences if stage $s$ is appended. We measure the resulting
degree imbalance as
%
\begin{equation}
D_s =
\max_i \left(d_i+\delta_i(\mathcal{M}_s)\right)
-
\min_i \left(d_i+\delta_i(\mathcal{M}_s)\right).
\label{eq:degree_spread}
\end{equation}
%
Stages with smaller $D_s$ are preferred. This criterion prevents repeated
selection of a small subset of predecessors and keeps access to the preceding
layer as uniform as possible.

Degree balance alone, however, cannot distinguish stages that use predecessors
equally often but repeatedly combine signals with similar structural history.
We therefore use the ancestry sets introduced in
Section~\ref{sec:representation} only as a secondary tie-breaking criterion.
For a candidate stage, we define its ancestry novelty as
%
\begin{equation}
N_s =
\frac{1}{|Q_s|}
\sum_{(a,b)\in Q_s}
\left(
1-
\frac{|A_a\cap A_b|}
{\max\!\left(1,\min(|A_a|,|A_b|)\right)}
\right),
\label{eq:ancestry_novelty}
\end{equation}
%
where $Q_s$ contains all pairs of the candidate matching, or at most 2,048
evenly sampled pairs for very large stages. Identical nonempty ancestry sets
contribute zero, whereas disjoint sets contribute one. Hence, among stages
with the same degree spread, LogicConnect favors the one that combines
predecessors with less redundant reachable support. The ancestry score does
not override degree balance; it is used only when the primary criterion
cannot distinguish between candidate stages.

Any remaining ties are resolved by a deterministic cyclic scale order, with
every nominal scale used once before the cycle restarts. Importantly,
LogicConnect selects \emph{whole predefined matching stages}. It does not
optimize individual edges, swap connections according to validation accuracy,
or perform an accuracy-driven search over candidate graphs. Only the final
truncated stage, required when the remaining number of gates is smaller than a
complete matching, is selected at the pair level. This stage-based
construction substantially reduces the connectivity search space while
retaining controlled predecessor reuse and multiscale signal mixing.

Algorithm~\ref{alg:construction} summarizes the hidden-layer construction.

\begin{algorithm}[t]
\caption{Degree-Balanced Multiscale Matching}
\label{alg:construction}
\begin{algorithmic}[1]
\REQUIRE Number of predecessors $n$, number of gates $m$, ancestry sets
$\{A_i\}$
\ENSURE Fixed predecessor pairs for the layer
\STATE Generate $K=\lceil\log_2 n\rceil$ multiscale candidate matchings
\STATE Initialize predecessor degrees $d_i\leftarrow 0$
\STATE Initialize cyclic scale order
\WHILE{unassigned gates remain}
    \STATE Determine whether a complete or partial stage is required
    \FOR{each eligible candidate stage $s$}
        \STATE Compute degree spread $D_s$
    \ENDFOR
    \STATE Retain stages with minimum $D_s$
    \IF{multiple stages remain}
        \STATE Compute ancestry novelty $N_s$ for the tied stages
        \STATE Retain stages with maximum $N_s$
    \ENDIF
    \STATE Resolve any remaining tie using the cyclic scale order
    \STATE Append the selected complete or partial matching
    \STATE Update predecessor degrees and ancestry
\ENDWHILE
\RETURN Fixed predecessor pairs
\end{algorithmic}
\end{algorithm}


\subsection{Convolutional DLGNs}
\label{sec:conv}

The same connectivity principles can be applied to convolutional DLGNs, but
the object being connected changes from individual signals to shared channel
groups. Convolutional DLGNs already impose a strong spatial inductive bias
through their receptive-field sampler and shared logic trees
\cite{petersen2024}. LogicConnect therefore does not alter the spatial
sampling pattern or replace the convolutional computation. Instead, it
constructs the \emph{channel pairs} presented to the shared logic trees,
allowing connectivity to be structured while preserving the original
convolutional execution model.

More specifically, the receptive-field sampler, tree depth, gate sharing,
and pooling operations remain identical to the baseline convolutional DLGN.
For the first convolutional layer, channel assignments are constructed using
the available input metadata, analogous to the semantic pairing described in
Section~\ref{sec:semantic}. Subsequent channel groups are formed using the
degree-balanced multiscale matching of
Section~\ref{sec:multiscale}. Since the same channel pairing is reused across
spatial positions, the construction preserves weight sharing and does not
introduce location-dependent routing. Thus, LogicConnect changes which
channels are combined, while the spatial locations sampled by the
convolutional operator remain unchanged.

To support the same structural criteria across convolutional layers, ancestry
is tracked at the channel level. Initially, each encoded input channel is
assigned a source label. After a shared logic tree produces a new output
channel, its ancestry is formed from the union of the ancestries of its input
channels. We additionally retain an identity label for each generated output
channel. This distinction is important when two channels obtain identical
upstream support: although their reachable sources are the same, they remain
separate signals and may learn different Boolean functions. As in the dense
case, these labels are used only during offline graph construction and do not
become trainable features.

After the convolutional blocks, the resulting feature maps are flattened for
classification. Each classifier input retains its channel ancestry together
with its spatial-feature identity, allowing the fixed classifier graph to be
constructed using the same source-aware pairing and multiscale matching
principles introduced above. Therefore, the convolutional implementation does
not require a separate connectivity-learning mechanism. LogicConnect modifies
only the fixed predecessor assignments; the convolutional sampler, shared
logic trees, pooling, gate parameterization, training objective, and
discretization procedure remain unchanged from the underlying convolutional
DLGN.

\begin{comment}
    


LogicConnect treats fixed DLGN connectivity as an offline architectural design problem and constructs the graph in two stages: source-aware pairing where signal semantics are known, followed by degree-balanced multiscale mixing where direct semantics disappear. It produces an ordinary fixed DLGN graph without learning routing or changing the declared architecture.


A. Design Principles and Overview --> This subsection should be new and short—approximately 2 paragraphs.
Something like "LogicConnect is entirely architecture-driven: the constructor does not
observe class labels, training loss, validation accuracy, or trained gate
parameters, and it performs no accuracy-driven edge selection or
individual-edge swapping."


B. Fixed-Graph Representation
Ancestry needs a more precise definition
A_j denotes reachable structural support, not functional dependence. A learned gate may ignore one of its accessible inputs.

C. Semantic Pairing at the Input
Why semantics?

The first-layer inputs correspond to known encoded image sources. Therefore, unlike hidden features, there is information about:

raw spatial position,
channel identity,
threshold/bit identity.

Then define exactly what constitutes a source. That makes the rule
“exclude pairs whose two bits come from the same raw source”
define the candidate families

Rather than prose alone, I would enumerate compactly:

C=Cspatial∪Cchannel∪Cencoded


D. Multiscale matching in hidden layers
This is the conceptual core of the deeper constructor.
Once signals have passed through the first logic layer, their indices no longer correspond directly to image coordinates; therefore, imposing the same semantic pairing rule at every depth would be unjustified.
We want:

balanced exposure,
systematic cross-signal mixing,
multiple interaction scales,
deterministic construction.

Then introduce the matchings.
A small \(n=8\) illustration would help tremendously:

E. Matching Selection: Balance First, Diversity Second
Separate the candidate matching definition from the matching-selection rule. Right now both are packed into one subsection.
Very important: stage-level selection: Importantly, LogicConnect selects predefined matching stages rather than
optimizing individual edges. Except for the final truncated stage required
to match the requested gate count, no edge-level search, replacement, or
accuracy-driven swapping is performed.

F. Applying the Same Constructor to Shared Convolutional Logic

\end{comment}